%% capitoli/cap3_legge_2026.tex
%% In testa al file, prima di scrivere, incollare l'estratto di
%% output/F4-indice.md relativo a questo capitolo (tesi sostenuta,
%% fonti, lunghezza stimata): serve a controllare gli scostamenti.
%% Convenzione delle etichette: sez:cap3:<nome>, tab:cap3:<nome>.

% Nucleo del tema: ogni affermazione dichiara su quale dei due piani si
% colloca (norma self-executing / criterio di delega). Una tabella
% sinottica delle due colonne e' ammessa: \label{tab:cap3:sinottica}

\chapter{La legge n. 1 del 2026: struttura e ratio dell'intervento}
\label{cap:cap3}

\section{Il disegno complessivo: responsabilità amministrativa, controlli, consulenza}
\label{sez:cap3:disegno}


\section{La consulenza come contrappeso alla ridefinizione della colpa grave}
\label{sez:cap3:contrappeso}


\section{L'art. 2: il passaggio dall'astratto al concreto}
\label{sez:cap3:art2}


\section{La delega: principi e criteri direttivi, termini, vincoli}
\label{sez:cap3:delega}


\section{Disposizioni immediatamente precettive e disposizioni che attendono il decreto legislativo}
\label{sez:cap3:precettivo}


